% !TEX root = ../tree_notes.tex
\chapter{Problem-Solving Strategies}\label{ch:app-strategies}

\section{Recognising a Tree Problem}

\begin{tcolorbox}[mathbase, colback=lightgray, colframe=primaryblue, title={Problem indicators}, colbacktitle=primaryblue]
\textbf{It is a tree.} Words such as \emph{node, root, parent, child, leaf, subtree,
sibling, ancestor, descendant, hierarchy}; ``$n$ nodes and $n-1$ edges''; ``connected
and acyclic''; an org chart, a file system, a family tree, a parse expression.

\textbf{Binary specifically.} Each node has a \code{left} and a \code{right}; a BST,
a heap, a segment tree; LeetCode \code{TreeNode*} input.

\textbf{Subtree queries.} ``For every node, the sum / max / count over its subtree'';
ask for a value that aggregates a whole branch. Root the tree and run one post-order
DFS.

\textbf{Path queries.} ``Between nodes $u$ and $v$'', ``the longest path'', ``maximum
path sum''. Think LCA and the identity
$\dist(u,v)=\depth(u)+\depth(v)-2\depth(\lca(u,v))$, or heavy--light decomposition.

\textbf{Rerooting.} ``Compute the answer for every node as if it were the root''.
Two DFS passes: gather from children, then push corrected values back down.
\end{tcolorbox}

\section{Constraints Suggest the Method}

\begin{center}
\renewcommand{\arraystretch}{1.2}
\begin{tabular}{@{}p{0.42\linewidth}p{0.5\linewidth}@{}}
\toprule
\textbf{Constraint / wording} & \textbf{Usually means} \\
\midrule
$n\le 2\cdot10^{5}$, $n-1$ edges & adjacency list, one DFS/BFS, $O(n)$ or $O(n\log n)$ \\
recursion depth can reach $n$ (a path) & write the traversal iteratively, or raise the stack \\
many queries on one static tree & preprocess once: depths, parents, Euler tour, binary lifting \\
``subtree of every node'' & post-order DFS, or Euler tour $+$ Fenwick/segment tree \\
``path between two nodes'' & LCA ($O(\log n)$ per query after $O(n\log n)$ preprocessing) \\
values on edges / weights & root the tree, keep the parent to orient each edge \\
``closest marked node'', ``distance to nearest'' & multi-source BFS or a rerooting DP \\
counting labeled trees / shapes & Cayley $n^{n-2}$, Catalan $C_n$ (Chapter 1) \\
tree is actually a BST in disguise & use the in-order $=$ sorted property (\chref{c05}) \\
\bottomrule
\end{tabular}
\end{center}

\section{A Five-Step Framework}

\begin{tcolorbox}[mathbase, colback=white, colframe=successgreen, title={From statement to accepted}, colbacktitle=successgreen]
\textbf{1. Root it.} Pick a root (node $1$, or the node the problem names). Almost
every tree algorithm is a DFS that carries the parent down and combines child answers
on the way up.

\textbf{2. Draw a small tree.} Six to ten nodes. Mark the depth, height and size of
each. Most tree recurrences become obvious once you can see the subtree of one node.

\textbf{3. State the recurrence.} Write what a node needs from its children (the DP
state) and what it returns to its parent. Name the base case: usually \nil{} or a leaf.

\textbf{4. Implement with a parent guard.} On an adjacency list, skip the parent. On
pointers, handle \nil{} first. Choose \code{long long} when a subtree sum or a count of
shapes can exceed $2^{31}$.

\textbf{5. Check the shapes.} The empty tree, a single node, a path (max depth), a
star (one root, many leaves), and a perfect tree. Off-by-one in depth vs height shows
up on exactly these.
\end{tcolorbox}
